% !TEX root = ../main.tex

%----------------------------------------------------------------------------------------
% APPENDIX A
%----------------------------------------------------------------------------------------

\chapter{Annexos}
\label{AppendixA}

\section{Estructura del document}

Aquest annex conté els documents complementaris al projecte: pressupostos, compliment normatiu i manual d'usuari.

%----------------------------------------------------------------------------------------

\section{Pressupost detallat}

El pressupost del projecte és de \textbf{0 euros} ja que totes les eines i serveis utilitzats són de codi obert i gratuïts:

\begin{table}[h]
\centering
\caption{Pressupost detallat del projecte}
\begin{tabular}{|l|l|c|}
\hline
\textbf{Concepte} & \textbf{Descripció} & \textbf{Cost} \\
\hline
\textbf{Desenvolupament} & 258 hores @ 0€/h (estudiant) & 0€ \\
\hline
Node.js 18 & Runtime JavaScript & 0€ \\
\hline
React 18 & Framework frontend & 0€ \\
\hline
Vite 5 & Bundler & 0€ \\
\hline
GitHub & Repositori i col·laboració & 0€ \\
\hline
GitHub Actions & CI/CD (2.000 min/mes gratuït) & 0€ \\
\hline
GitHub Pages & Desplegament frontend & 0€ \\
\hline
VS Code & Editor de codi & 0€ \\
\hline
LaTeX & Preparació documentació & 0€ \\
\hline
Draw.io & Diagrames & 0€ \\
\hline
\textbf{TOTAL} & & \textbf{0€} \\
\hline
\end{tabular}
\end{table}

\textbf{Inversió de temps:}

\begin{table}[h]
\centering
\caption{Distribució del temps de desenvolupament}
\begin{tabular}{|l|c|}
\hline
\textbf{Fase} & \textbf{Hores} \\
\hline
Investigació i configuració & 30 \\
\hline
Desenvolupament ETL & 60 \\
\hline
Frontend (Dashboards 1-2) & 67 \\
\hline
Frontend (Dashboards 3-4) & 49 \\
\hline
Automatització i documentació & 52 \\
\hline
\textbf{TOTAL} & \textbf{258} \\
\hline
\end{tabular}
\end{table}

%----------------------------------------------------------------------------------------

\section{Compliment normatiu}

\subsection{Protecció de dades (LOPDGDD)}

El compliment del Reglament General de Protecció de Dades (RGPD) i la Llei Orgànica 3/2018 de Protecció de Dades Personals (LOPDGDD) es garanteix pels següents aspectes:

\begin{itemize}
    \item \textbf{Dades públiques}: Totes les dades utilitzades pel sistema provenen de fonts públiques de la Generalitat de Catalunya, publicades sota llicències obertes. No es tracten dades de caràcter personal.
    \item \textbf{No recollida de dades d'usuaris}: El frontend no recull, emmagatzema ni processa cap informació personal dels usuaris. No hi ha formularis de registre, sistemes de login ni cookies de seguiment.
    \item \textbf{Transparència}: El codi font és completament obert, permetent que qualsevol persona verifiqui què fa el sistema amb les dades.
\end{itemize}

\subsection{LSSICE}

La Llei 34/2002 de Serveis de la Societat de la Informació i de Comerç Electrònic (LSSICE) estableix requisits d'informació per als serveis en línia. En tractar-se d'un projecte acadèmic sense finalitat comercial, els requisits es redueixen, però s'ha inclòs:

\begin{itemize}
    \item Informació clara sobre l'autoria del projecte a la pàgina \texttt{/about}.
    \item Identificació de les fonts de dades utilitzades (APIs Socrata de la Generalitat).
    \item Enllaç al repositori de codi font per a verificació independent.
\end{itemize}

\subsection{Accessibilitat}

L'accessibilitat del sistema s'ha abordat seguint les pautes WCAG 2.1 de forma progressiva:

\begin{itemize}
    \item \textbf{Contrast adequat}: S'han seleccionat colors amb contrast suficient entre text i fons per garantir la llegibilitat.
    \item \textbf{Etiquetes descriptives}: Tots els elements interactius tenen etiquetes clares i comprensibles.
    \item \textbf{Navegació per teclat}: Els elements interactius principals són accessibles mitjançant teclat.
    \item \textbf{Disseny responsiu}: L'aplicació s'adapta a diferents mides de pantalla, facilitant l'accés des de dispositius mòbils.
\end{itemize}

Tot i aquests esforços, es reconeix que algunes visualitzacions complexes (com els gràfics de Recharts) podrien millorar la seva accessibilitat amb text alternatiu més detallat i suport per a lectors de pantalla, aspecte que es podria abordar en futures iteracions.

%----------------------------------------------------------------------------------------